\documentclass[11pt]{article}
\usepackage[margin=1in]{geometry}
\usepackage{amsmath,amssymb,amsthm}
\usepackage{booktabs}
\usepackage{hyperref}
\usepackage{xcolor}
\usepackage{siunitx}

\title{Independent Replication --- OSTI 2564727 \\
\large Accurate Numerical Simulations of Open Quantum Systems Using Spectral Tensor Trains}
\author{Independent replication (OSTI-100 set, applied\_math rank 5)}
\date{2026-07-02}

\begin{document}
\maketitle

\begin{abstract}
We independently replicate the theoretical anchor of Grimm \& Eaves,
\emph{Accurate Numerical Simulations of Open Quantum Systems Using Spectral Tensor Trains}
(J.\ Chem.\ Phys.\ 2025, DOI \texttt{10.1063/5.0228873}, OSTI 2564727).
Q-ASPEN (Quantum Accelerated Stochastic Propagator Evaluation) solves the
Stochastic Liouville Equation (SLE) by representing a Gaussian influence
kernel as a spectral tensor train (STT).  Its load-bearing exact statement is that
in the extrinsic Markov limit the noise-averaged SLE reduces \emph{exactly}
to the Lindblad master equation (Eqs.~16--17).  Two mutually independent routes
--- direct RK4 integration of the Lindblad ME, and Monte-Carlo averaging of
SLE stochastic-Schr\"odinger trajectories --- agree to Monte-Carlo sampling
precision ($\sim 10^{-3}$ at $N_{\mathrm{traj}}{=}40\,000$) across three
physically distinct qubit systems, with a closed-form dephasing case matched to
$9.6\times 10^{-13}$.  \textbf{Verdict: REPLICATED} for the core analytic claim;
PARTIAL for the paper overall (STT machinery, memory-scaling claim, and
intrinsic-noise spin-boson comparison were not re-run).
\end{abstract}

\section{Paper summary}

Grimm and Eaves introduce Q-ASPEN, a numerically-exact method for the
time-dependent noise-averaged reduced density matrix
$\langle \rho(t)\rangle$ of an open quantum system driven by prescribed colored
\emph{intrinsic} (thermal/quantum, obeying fluctuation--dissipation) and
\emph{extrinsic} (external-field) noise.  The Stochastic Liouville Equation
$d\rho/dt = -i L(t)\rho$ with $L(t)=L_0+\xi(t)L_1$ is solved by
(i)~Trotter-discretizing the noise-averaged propagator $\Phi_N$ (Eqs.~4--6),
(ii)~casting the Gaussian influence kernel $K$ (Eqs.~7--8) as a Boltzmann-like
weight over eigenfrequency walks (a ``statistical mechanics of trajectories''
picture), and (iii)~representing $K$ as a spectral tensor train
of Chebyshev-expanded matrix-valued components $K_\alpha(\omega)$
(Eqs.~10--14), trained by SGD/importance sampling (which they observe to be
obstructed by a barren-plateau phenomenon).  Benchmarks include a spin-boson
model (Fig.~3, matching PT-TEMPO) and a 2--32 site extrinsic-noise quantum
chain (Fig.~4; reported memory scaling $p\!\approx\!2.0$ for Q-ASPEN vs
$p\!\approx\!8.3$ for PT-TEMPO).

\paragraph{Load-bearing exact result (Eqs.~15--17).}
In the extrinsic Markov limit $H=H_0+\xi(t)V$ with
$\langle \xi(t)\xi(s)\rangle=\gamma\delta(t-s)$, the correlation matrix collapses
to $G=\gamma\tau\mathbf{1}$ and the fugacities become
$Z_\alpha = e^{-iL_0\tau/2}\, e^{-\gamma L_1^2 \tau/2}\, e^{-iL_0\tau/2}$.
Applying Baker--Campbell--Hausdorff and taking $\tau\to 0$ yields
$\Phi(t)=\exp\!\big({-}itL_0 - t\gamma L_1^2/2\big)$, equivalent to the
Lindblad master equation
\begin{equation}
\frac{d\langle \rho\rangle}{dt}
= -i\,[H_0,\langle \rho\rangle]
+ \gamma\!\left(V\langle \rho\rangle V - \tfrac{1}{2}\{V^2,\langle \rho\rangle\}\right).
\label{eq:lindblad}
\end{equation}
This is the analytic anchor on which the correctness of Q-ASPEN rests, and the
principal target of the present replication.

\section{Claims table}

\begin{center}
\begin{tabular}{clllll}
\toprule
ID & Claim & Type & Independently testable? & Tested here? \\
\midrule
C1 & Extrinsic-Markov SLE $\Rightarrow$ exact Lindblad ME & analytic & Yes & \textbf{Yes} \\
C2 & Method conserves trace / gives physical $\rho$ & numerical & Yes (via C1) & Yes ($\operatorname{Tr}{=}1.0000$) \\
C3 & Spin-boson (Fig.~3) matches PT-TEMPO & numerical & Partially & Geometry only \\
C4 & 32-site chain memory scaling $p\!\approx\!2.0$ vs $8.3$ & scaling & Full STT+OQuPy port & No (out of scope) \\
C5 & STT training exhibits barren plateaus & empirical & Full STT training & No \\
\bottomrule
\end{tabular}
\end{center}

\section{Method (independent, from equations only)}

\textbf{Reference (exact Lindblad).}  RK4 integration of
Eq.~\eqref{eq:lindblad} (\texttt{sle\_lindblad.py:integrate\_lindblad}).
Validated against the closed form: for $V=\sigma_z$, $H_0=0$,
Eq.~\eqref{eq:lindblad} gives $\rho_{01}(t)=\rho_{01}(0)e^{-2\gamma t}$.

\textbf{Independent stochastic route.}  We Monte-Carlo average the SLE
stochastic Schr\"odinger equation
$i\,d|\psi\rangle/dt = (H_0 + \xi(t)V)|\psi\rangle$ with real white noise.
Each step draws $\Delta_k \sim \mathcal{N}(0,\gamma\tau)$ and applies the
symmetric Trotter step $e^{-iH_0\tau/2}\,e^{-i\Delta_k V}\,e^{-iH_0\tau/2}$
(mirroring Eq.~5), then averages $\rho=|\psi\rangle\langle\psi|$ over
trajectories.  For real noise coupling to Hermitian $V$, each trajectory is
unitary; decoherence emerges purely from ensemble averaging --- the Kubo
stochastic-Liouville mechanism the paper builds on.  Vectorized across
trajectories for speed (exact reformulation, not an approximation).

\textbf{Systems tested.}
(1) Pure dephasing: $V=\sigma_z$, $H_0=(\varepsilon/2)\sigma_z$;
(2) Transverse coupling: $V=\sigma_x$, $H_0=(\varepsilon/2)\sigma_z$
(non-commuting);
(3) Biased qubit (Fig.~3 geometry):
$H_0=\Omega\sigma_x + \varepsilon\sigma_z$, $V=\alpha\sigma_z$ with
$\Omega=1$, $\varepsilon=0.5$, $\alpha=0.75$.

\textbf{Stack.}  Python 3.14 + \texttt{numpy} + \texttt{matplotlib}.
Commands: \verb|python3 sle_lindblad.py|, \verb|python3 convergence.py|,
\verb|python3 plot_dynamics.py|.

\section{Results vs paper}

\paragraph{Analytic anchor.} For $V=\sigma_z$, $H_0=0$, $\gamma=0.5$:
Lindblad $\rho_{01}(t)$ vs $0.5\,e^{-2\gamma t}$ agrees to
\textbf{max error $9.6\times 10^{-13}$} (machine precision), confirming
correct conventions.

\paragraph{C1 --- SLE Monte-Carlo vs exact Lindblad}
($N_{\mathrm{traj}}=40\,000$, $\tau=0.01$, $t\in[0,4]$):

\begin{center}
\begin{tabular}{lccl}
\toprule
Case & $\max|\rho_{\mathrm{MC}}-\rho_{\mathrm{Lindblad}}|$ &
$\operatorname{Tr}(\rho_{\mathrm{MC}})$ & $\langle \sigma_z\rangle$ tracking \\
\midrule
1. Pure dephasing $V=\sigma_z$   & $4.3\times 10^{-3}$ & 1.0000 & flat (correct) \\
2. Transverse $V=\sigma_x$       & $4.2\times 10^{-3}$ & 1.0000 & $1\to 0$, matches \\
3. Biased qubit (Fig.~3 geom.)   & $2.9\times 10^{-3}$ & 1.0000 & oscillates, 3-dec.\ match \\
\bottomrule
\end{tabular}
\end{center}

Residuals lie at the Monte-Carlo sampling floor ($\sim 1/\!\sqrt{N}\approx
5\times 10^{-3}$ at $N=40\,000$).

\paragraph{Convergence (no systematic bias).}
Errors decrease with trajectory count:
$7.5\!\times\!10^{-3}$ (2.5k) $\to$ $5.4\!\times\!10^{-3}$ (10k) $\to$
$2.8\!\times\!10^{-3}$ (40k) $\to$ $2.6\!\times\!10^{-3}$ (160k).
A Trotter-step scan confirms the residual floor is sampling + $\mathcal{O}(\tau^2)$
discretization, not model mismatch.  As $N\to\infty$, $\tau\to 0$, the SLE
ensemble average approaches the exact Lindblad solution, exactly as
Eqs.~16--17 claim.

\paragraph{Not reproduced (honest scope).}  The STT/tensor-train machinery,
its $\mathcal{O}(d^2)$ memory-scaling claim (C4), the intrinsic-noise
spin-boson vs.\ PT-TEMPO comparison (C3), and barren-plateau training (C5)
were out of scope; they require porting the full STT method and (for C3)
the App.~A kernel plus OQuPy.

\section{Independent LLM judge (free Argo \texttt{gpt-5.2})}

Verdict: \textbf{REPLICATED} for the core analytic claim.  Verbatim gist:
``Independent Monte-Carlo averaging of the SLE converges toward the
Lindblad master-equation solution across several distinct qubit
Hamiltonians/couplings, with an analytic dephasing case matched to machine
precision and no evidence of systematic deviation.''
Full text: \texttt{evidence/llm\_judge\_verdict.txt}.

\section{GENUINE CRITIQUE}

The following are honest reservations that a hurried ``verdict = REPLICATED''
summary should not obscure.

\begin{enumerate}
\item \textbf{We replicated the anchor, not the machine.}  The paper's
selling point is the STT compression and its $p\!\approx\!2.0$ vs $8.3$
memory scaling (Fig.~4).  That claim is what would justify replacing
PT-TEMPO in practice, and it is precisely what we did \emph{not} re-derive.
Our replication validates that the SLE $\to$ Lindblad limit is exact ---
which is a well-known Kubo/stochastic-Liouville result and, honestly,
would have been surprising if it had \emph{failed}.  Passing it is
necessary but not sufficient for endorsing Q-ASPEN.

\item \textbf{Extrinsic Markov is the easiest corner of the space.} White
noise + real coupling to a Hermitian operator gives unitary per-trajectory
dynamics and no memory kernel.  The hard cases in the paper are colored
intrinsic noise (App.~A kernel obeying FDT) and the non-Markovian
spin-boson benchmark of Fig.~3; those require the full STT plus a
PT-TEMPO reference and were not re-run.  The residual $\sim\!3\times 10^{-3}$
error we report is genuinely at the sampling floor, but only for the easy
regime.

\item \textbf{No PT-TEMPO cross-check.}  A stronger replication would run
OQuPy/PT-TEMPO on the exact Fig.~3 geometry (biased qubit + bosonic bath,
spectral density and $\beta$ as published) and overlay both curves.  We
mimicked the geometry but drove it with white extrinsic noise, so our
``matches Fig.~3'' language means \emph{our SLE-MC matches our own
Lindblad}, not that either matches the published spin-boson dynamics.

\item \textbf{Barren-plateau claim (C5) is a training pathology, not a
correctness bug.}  It is important for whether Q-ASPEN is \emph{usable}
in practice, but it is orthogonal to the analytic anchor.  A future
replication should tabulate the empirical distribution of gradient norms
across STT layers on a small system and confirm the plateau structure,
which we did not attempt.

\item \textbf{Compute honesty.}  This replication used local NumPy on a
laptop-class machine.  We have no independent handle on the wall-clock or
memory claims of the paper, only on the mathematics.  The $p\!\approx\!2.0$
scaling claim is untested by us.

\item \textbf{Numerical caveats.}  Our biased-qubit case matches to
$\sim\!3\times 10^{-3}$ ``to 3 decimals,'' which is $N_{\mathrm{traj}}^{-1/2}$
noise; sub-percent claims about oscillation phase would require
$N_{\mathrm{traj}}\gtrsim 10^{6}$ or a variance-reduction scheme (antithetic
noise, control variates).  We did not do this.
\end{enumerate}

\textbf{Net position.}  The theoretical foundation of Q-ASPEN is solid and
reproduces independently.  The engineering claims (STT compression, memory
scaling, training behavior, non-Markovian benchmarks) remain open in this
replication and deserve a follow-up with the full method + OQuPy.

\section{Assessment and verdict}

The paper's central exact statement --- that noise-averaging the SLE with
Markovian white noise yields precisely the Lindblad master equation
(Eqs.~16--17) --- is independently reproduced.  Two mutually independent
routes agree to sampling + Trotter precision on three physically distinct
qubit systems; a closed-form dephasing case matches to $10^{-12}$.  This
validates the theoretical foundation of Q-ASPEN.  The STT approximation and
scaling claims were not re-derived and are marked untested.

\medskip
\noindent\textbf{Verdict:} \textcolor{blue}{\textbf{REPLICATED}} for the core
analytic claim (Eqs.~16--17);
\textbf{PARTIAL} for the paper overall (STT machinery and Fig.~4 scaling
benchmarks not re-run).

\end{document}
